% !PRINTED 181
\weberpage{181}

%% [running head, not body text -- commented out 2026-08-07] Page furniture (running
%% head, folio, signature) is excluded from this edition by standing policy.
%% \noindent \S.\,47. \hfill \textit{Normalformen.} \hfill 181

But we can also conclude conversely that, if $S$ possesses a normal divisor of index $e$, the form problem is imprimitive. For in order to obtain a system of the functions $X_1, \ldots X_m$, we need only take an invariant $\varrho$ of the group $T$ which contains $e$ different values in $S$, and we can take for $X_1, \ldots X_m$ simply the values $\varrho, \varrho_1, \ldots \varrho_{e-1}$. The group $\Sigma$ is then the permutation group called forth by $S$ among the $\varrho$.

In the sense of \S.\,44, however, it will always be a matter of making $m$ as small as possible, and a reduction of the problem in this sense will be achieved only when $m < n$.

\begin{center}
\S.\,47.

\gesperrt{Normal Forms.}
\end{center}

We must undertake one more general consideration before we pass on to special applications. We have already seen in \S.\,37 that from every group $S$ of linear substitutions we can derive infinitely many isomorphic groups, by transforming with an arbitrary substitution $L$ of the same dimension, that is, by forming the transformed group
\weq{1}{L^{-1} S L}
and this is equivalent to the introduction of other variables $y$ in place of $x$ by the substitution
\weq{2}{(x) = L(y).}

We can use this transformation to represent the group $S$ in a simple normal form, and then the invariants also receive certain fixed normal forms, which in many cases can assume very simple shapes. If we form for the normal form the resolvent of the form problem [\S.\,43, (8)], which we now wish to write in the form:
\weq{3}{\Phi(\Theta, A_1, A_2, \ldots) = 0,}
wherein $A_1, A_2, \ldots$ denote invariants of the group $S$, then this too, if we use the normal form, will receive a simple shape. The variables $x_i$ can, as we have seen
%CONT>>
